\section{The remaining quantum recurrence}
\label{ct22-next}

The actual stationary substitution is now represented by
strictly supported functions of the next-coarse variables.
The connected activities supplied by Notebook 21, and its
separate normalization terms after a common reduction of
their support exponent, satisfy the input form of the
transport theorem. Thus their sum can be expressed as
local next-coarse activities on the explicitly smaller
analytic ball. This is a completed finite-step interface.

For the specific convention used in Notebook 21, the
integration cells are already the next-coarse block cells.
That substitution uses the identity projection. Projection
into another set of $L$-blocks occurs only when preparing
these activities for the following integration. It cannot
be counted twice.

A repeated quantum construction requires more. The generated
activity norm has a fixed small support exponent, whereas the
Gaussian integration theorem asks for a stronger weighted
input. The connected-carrier estimate explains a possible
geometric gain at the later projection, but its constants
and the new fluctuation-chart enlargement must be included.
Notebook 24 makes that inherited-interaction budget explicit;
Notebook 26 meets it using a fixed auxiliary grouping and
recovers a common positive carrier exponent.
The smaller analytic radius in the transport theorem must
also be replaced by a proved stable domain recursion.
The explicit pure-gauge example above rules out obtaining
that recursion merely by declaring that every link angle
contracts by $L^{-1}$ in this rooted representative.
Notebook 25 constructs a smaller pre-root-gauge representative
and identifies the analytic composition germ for an invariant
total functional, with its remaining activity and integral
domain requirements kept explicit.

Finally, the vacuum, marginal gauge-action coefficient and
irrelevant interactions need their appropriate extractions
and scaling estimates. A bounded norm map is insufficient:
vacuum singleton terms already grow by block volume.
Large-field complements, source insertions and the continuum
reconstruction requirements of Notebook 01 remain part of
the full existence argument.

\begin{thebibliography}{9}
\bibitem{prior22}
\emph{Toward a four-dimensional Yang--Mills construction},
research Notebooks 15--21 in this workspace. The physical
stationary branch, exact current pivot chart, retained
inverse, scalar support construction and connected
integration activities are the principal inputs.
\bibitem{balaban10922}
T.~Ba\l aban,
\emph{Renormalization group approach to lattice gauge field
theories I}, Commun. Math. Phys. \textbf{109} (1987), 249--301,
\href{https://scispace.com/pdf/renormalization-group-approach-to-lattice-gauge-field-3ygai0vxt1.pdf}{original article copy}.
\bibitem{balaban11622}
T.~Ba\l aban,
\emph{Renormalization group approach to lattice gauge field
theories. II. Cluster expansions},
Commun. Math. Phys. \textbf{116} (1988), 1--22,
\href{https://scispace.com/pdf/renormalization-group-approach-to-lattice-gauge-field-4penm4k7hc.pdf}{original article copy}.
\bibitem{source22}
E.~Sussman and collaborators,
\href{https://github.com/ethanwsussman/Yang--Mills-LLM-experiment}{Yang--Mills-LLM-experiment},
deposited raw accounts at commit
\texttt{e392ef45013a581664bab76b90867170b30fe10e}.
\end{thebibliography}
\end{document}
